\documentclass[10pt,a4paper]{amsart}
\usepackage[margin=19mm]{geometry}
\usepackage{amsmath,amssymb,mathtools}
\usepackage[scr=rsfs,cal=euler]{mathalfa}
\usepackage{xfrac}
\usepackage[unicode,hidelinks,bookmarks=false]{hyperref}
\newcommand{\ie}{i.e.\ }
\newcommand{\cf}{cf.\ }
\newcommand{\resp}{respectively\ }
\newcommand{\Subsection}{\subsection}
\providecommand{\nobreakdash}{\nobreak\hbox{-}}
\setlength{\emergencystretch}{2em}
\multlinegap0pt
\usepackage{amssymb}
\usepackage{enumerate}
\usepackage{float}
\usepackage{xcolor}
\usepackage{tcolorbox}
\newtheorem{lemma}{Lemma}
\newtheorem{proposition}{Proposition}
\title{Cullen and Woodall numbers in Padovan and Perrin sequences}
\author{Herbert Batte, Eric F. Bravo, Florian Luca}
\date{Source: arXiv:2605.23084v1 (CC BY 4.0)}
\begin{document}
\maketitle

\noindent\emph{Source and licence.} Cullen and Woodall numbers in Padovan and Perrin sequences. Version arXiv:2605.23084v1 (CC BY 4.0), distributed by arXiv under the Creative Commons Attribution 4.0 International licence, http://creativecommons.org/licenses/by/4.0/, as recorded in the arXiv licence metadata for this exact version and shown on the abstract page https://arxiv.org/abs/2605.23084v1. Full licence notice: \texttt{licenses/arxiv-2605.23084-CC-BY-4.0.txt}.\par\smallskip

\noindent\textit{Editorial scope.} Counted unit: the article's proposition prop:Pn (source0.tex lines 484--499). The article's complete original proof of it is delivered with it (source0.tex lines 501--765, one \texttt{proof} environment of 265 lines). No mathematics is written, repaired or re-derived: the statement and the proof are the article's own lines, in the article's own notation, and no retained line is substituted or re-flowed.  Retained uncounted as the source-local context the proof needs: lines 171--177; lines 247--253; lines 256--267; lines 282--290; lines 306--320; lines 378--394; lines 1363--1363.  Nothing was omitted from any retained span, so the package carries no omission record.  No retained line contains a citation, so the wrapper carries no bibliography and no bibliography entry.  No retained line contains a figure environment, an image-inclusion command or drawing code, so no image is delivered and the package declares no asset dependency.  Editorial formatting only, in addition: an AMS \texttt{amsart} preamble that restates the article's own declarations and macros used by the retained lines, namely the environments \texttt{lemma}, \texttt{proposition} on separate counters, together with the standard packages those lines need.  The retained environments print in this document as Lemma 1, Proposition 1 in order of appearance, whereas the article prints the same items with the numbering of its own full document.  The complete native source of the article as distributed in the arXiv e-print for this exact version is bundled unchanged as \texttt{sources/201180/source0.tex}.
\begin{lemma}\label{eq:binet_padovan}
For all $n \in \mathbb{Z}$, we have
\begin{equation*}
	P_n = c_{\alpha}\alpha^n + c_{\beta}\beta^n + c_{\gamma}\gamma^n,
\end{equation*}
where $c_{x}:=(x^{2}+x)/(3x^{2}-1)$ for all $x\in \{\alpha,\beta, \gamma\}$.  
\end{lemma}

\begin{lemma}\label{L1}
	For all positive integers $r,s$ we have the relations:
	\begin{align*}
		P_{r+s}&=P_{s-1}P_{r-1}+P_{s-2}P_{r}+P_{s-3}P_{r-2},\\
		R_{r+s}&=P_{s-1}R_{r-1}+P_{s}R_{r-2}+P_{s-2}R_{r-3}.
	\end{align*}
\end{lemma}

\begin{lemma}\label{L2}
	For $n\ge 0$, we have 
	\begin{align*}
		\nu_2\left(P_{n}\right)= 
		\begin{cases}
			0, & \mbox{if } n\equiv 0,1,2,5,\pmod{7};\\
			\nu_2(n+4)+1, & \mbox{if } n\equiv 3\pmod{7};\\
			\nu_2((n+3)(n+17))+1, & \mbox{if }  n\equiv 4\pmod{7};\\
			\nu_2((n+1)(n+8))+1, & \mbox{if }  n\equiv 6\pmod{7}.
		\end{cases}
	\end{align*}
\end{lemma}

\begin{lemma}\label{L4}
	For the integers $j$ and $k,t\ge 1$, we get
	$$
	P_{7\cdot 2^{t}k+j}\equiv\begin{cases}
		P_{j}\pmod{2^{t+2}}, & \text{if } j\equiv 1,2,4\pmod{7},\\
		P_{j}+k\cdot 2^{t+1}\pmod{2^{t+2}}, & \text{otherwise}.
	\end{cases}
	$$
\end{lemma}

\begin{lemma}\label{lem:doubling}
	The following hold.
	\begin{itemize}
		\item[\rm(i)] $R_7 = 7$, $R_{14} = 51$, $R_{28} = 2627$, 
		$R_{-7} = -1$, $R_{-14} = -13$, $R_{-28} = 67$.
		\item[\rm(ii)] For all $k \ge 2$,
		\begin{equation*}
			R_{7\cdot 2^{k+1}} = R_{7\cdot 2^k}^2 - 2R_{-7\cdot 2^k}.
		\end{equation*}
		\item[\rm(iii)] For all $k \ge 2$,
		\begin{equation*}
			R_{-7\cdot 2^{k+1}} = R_{-7\cdot 2^k}^2 - 2R_{7\cdot 2^k}.
		\end{equation*}
	\end{itemize}
\end{lemma}

\begin{lemma}\label{lem:shift}
	Let $\{U_n\}_{n\in\mathbb{Z}}$ be any sequence satisfying 
	the recurrence $U_{n+3} = U_{n+1} + U_n$, and let $j$ be 
	any integer. Then for all $k \ge 1$,
	\begin{itemize}
		\item[\rm(i)]
		\begin{equation*}
			U_{j-7\cdot 2^{k+1}} = -U_j R_{7\cdot 2^k} 
			+ U_{j+7\cdot 2^k} + U_{j-7\cdot 2^k}R_{-7\cdot 2^k},
		\end{equation*}
		\item[\rm(ii)]
		\begin{equation*}
			U_{j+7\cdot 2^{k+1}} = -U_j R_{-7\cdot 2^k} 
			+ U_{j-7\cdot 2^k} + U_{j+7\cdot 2^k}R_{7\cdot 2^k}.
		\end{equation*}
	\end{itemize}
\end{lemma}

 \subsection{Py Code I}\label{app1}

\begin{proposition}\label{prop:Pn}
	For $n\ge 0$ such that $n\not\equiv 50\pmod{112}$, we have
	\begin{equation*}
		\nu_{2}(P_{n}+1)= 
		\begin{cases}
			0, & \text{if } n\equiv 3,4,6\pmod{7},\\
		\nu_{2}(n+9)+1, & \text{if } n\equiv 5\pmod{7},\\
			1, & \text{if } n\equiv 0,1.2,7,9\pmod{14},\\
			3, & \text{if } n\equiv 8\pmod{28},\\
			5, & \text{if } n\equiv 22\pmod{56},\\
			\nu_{2}(n+6)+4, & \text{if } n\equiv 106\pmod{112}.
		\end{cases}
	\end{equation*}
	Moreover, if $n\equiv 50\pmod{112}$ and $n<2.3\times 10^{15}$, 
	then $\nu_{2}(P_{n}+1)\le 56$.
\end{proposition}

\begin{proof}
	First, we verify that the seven listed residue classes are 
	mutually exclusive and exhaustive. Every integer $n\ge 0$ 
	satisfies exactly one of the relations $n\equiv 0,1,2,3,4,5,6\pmod{7}$. 
	The class $n\equiv 1\pmod{7}$ is further refined by the 
	remainder of $n$ modulo higher powers of $2$ by $7$, 
	which yields the subcases modulo $14$, $28$, $56$, and $112$. 
	It can be directly verified that the union of the seven classes 
	covers all non-negative integers exactly once.
	
	Now let's look at each case separately.
	
	\medskip
	\noindent\textbf{Case 1.} $n\equiv 3,4,6\pmod{7}$.
	
	By Lemma \ref{L2}, $\nu_2(P_n)\ge 1$ for $n$ in each of 
	these residue classes, so $P_n$ is even. Therefore, $P_n+1$ 
	is odd, which implies that $\nu_2(P_n+1)=0$.
	
		\medskip
	\noindent\textbf{Case 2.} $n\equiv 5\pmod{7}$.
	
	Let $n=7s+5$ for some integer $s\ge 0$. Then, $n+9=7(s+2)$. By the fundamental theorem of arithmetic, $s+2$ can be uniquely expressed as $2^{t}k$ for some integer $t\ge 0$ and some odd $k$. Therefore, every $n$ in this residue 
	class can be uniquely expressed as $n=7\cdot 2^t k-9$ for 
	some $t\ge 0$ and some odd $k$. Since 
	$j=-9\equiv 5\pmod{7}$, it follows from Lemma \ref{L4} that
	\begin{equation*}
		P_{7\cdot 2^t k-9}\equiv P_{-9}+k\cdot 2^{t+1}
		\pmod{2^{t+2}}.
	\end{equation*}
	It follows from Lemma \ref{eq:binet_padovan} that $P_{-9}=c_{\alpha}\alpha^{-9}+c_{\beta}\beta^{-9}+c_{\gamma}\gamma^{-9}=-1$. Therefore,
	\begin{equation*}
		P_n+1\equiv k\cdot 2^{t+1}\pmod{2^{t+2}}.
	\end{equation*}
	Since $k$ is odd, it follows that $\nu_2(P_n+1)=t+1$. 
	Finally, since $n+9=7\cdot 2^t k$ and $k$ is odd, we 
	have that $\nu_2(n+9)=t$, so $\nu_2(P_n+1)=\nu_2(n+9)+1$, as stated.
	
	\medskip
	\noindent\textbf{Case 3.} $n\equiv 0,1,2,7,9\pmod{14}$.
	
	We prove by weak induction that $P_n\equiv 1\pmod{4}$ for all such $n$, which implies that $P_{n}+1=4t+2$ for some integer $t\ge 0$ and so $\nu_2(P_n+1)=\nu_{2}(2(2t+1))=1$.
	
	\smallskip
	\textit{Subcase} $n\equiv 0,7\pmod{14}$. Let $n=7k$ for some integer 
	$k\ge 0$. The base case holds, since $P_0=1\equiv 1\pmod{4}$. 
	Suppose that $P_{7k}\equiv 1\pmod{4}$ for some fixed integer
	$k\ge 0$. Then, applying Lemma \ref{L1} with $r=7k+2$ and $s=5$, we obtain
	\begin{align*}
		P_{7(k+1)} = P_{(7k+2)+5} 
		&= P_{4}P_{7k+1} + P_{3}P_{7k+2} 
		+ P_{2}P_{7k}\\
		&= 2P_{7k+1}+2P_{7k+2}+P_{7k}\\
		&\equiv 2P_{7k+4}+1\pmod{4},
	\end{align*}
	where we used the fact that $P_2=1$ and $P_3=P_4=2$, and the recurrence relation of $\{P_{n}\}_{n}$. Since $7k+4\equiv 4\pmod{7}$, 
	Lemma~\ref{L2} states that $P_{7k+4}$ is even. Hence, $2P_{7k+4}\equiv 0\pmod{4}$. Therefore,
	\begin{equation*}
		P_{7(k+1)}\equiv 1\pmod{4}.
	\end{equation*}
	
	\smallskip
	\textit{Subcase} $n\equiv 2,9\pmod{14}$. Let $n=7k+2$ for some integer 
	$k\ge 0$. Since $P_2=1\equiv 1\pmod{4}$, the case base holds. 
	Suppose that $P_{7k+2}\equiv 1\pmod{4}$ for some fixed integer $k\ge 0$. Then, 
	applying Lemma~\ref{L1} with $r=7k+2$ and $s=7$, and 
	using the fact that $P_4=2$, $P_5=3$, and $P_6=4$, we obtain
	\begin{align*}
		P_{7(k+1)+2} = P_{(7k+2)+7} 
		&= P_{6}P_{7k+1}+P_{5}P_{7k+2}
		+P_{4}P_{7k}\\
		&= 4P_{7k+1}+3P_{7k+2}+2P_{7k}\\
		&\equiv 3+2P_{7k}\pmod{4}.
	\end{align*}
	Since $7k\equiv 0\pmod{7}$, Lemma~\ref{L2} states that $P_{7k}$ 
	is odd, so $2P_{7k}\equiv 2\pmod{4}$. Therefore,
	\begin{equation*}
		P_{7(k+1)+2}\equiv 1\pmod{4}.
	\end{equation*}
	
	\smallskip
	\textit{Subcase} $n\equiv 1\pmod{14}$. Let $n=14k+1$ for some integer $k\ge 0$. Since $P_1=1\equiv 1\pmod{4}$, the case base is satisfied. Suppose that $P_{14k+1}\equiv 1\pmod{4}$ for some fixed integer $k\ge 0$. Applying 
	Lemma~\ref{L1} with $r=14k+2$ and $s=13$, and using the fact that 
	$P_{10}=12$, $P_{11}=16$, and $P_{12}=21$, we obtain
	\begin{align*}
		P_{14(k+1)+1} = P_{(14k+2)+13}
		&= P_{12}P_{14k+1}+P_{11}P_{14k+2}
		+P_{10}P_{14k}\\
		&= 21P_{14k+1}+16P_{14k+2}+12P_{14k}\\
		&\equiv P_{14k+1}\equiv 1\pmod{4}.
	\end{align*}
	
	\medskip
	\noindent\textbf{Case 4.} $n\equiv 8\pmod{28}$.
	
	We will prove by weak induction that $P_n\equiv 7\pmod{16}$, which gives $P_{n}+1=16t+8$ for some integer $t\ge 0$, so 
	$\nu_2(P_n+1)=\nu_{2}(8(2t+1))=3$. Let $n=28k+8$ for an integer $k\ge 0$. Since $P_8=7\equiv 7\pmod{16}$, the case base holds. Suppose that $P_{28k+8}\equiv 7\pmod{16}$ for some fixed $k\ge 0$. Applying 
	Lemma~\ref{L1} with $r=28k+10$ and $s=26$, and using the fact that
	$P_{23}=465\equiv 1\pmod{16}$, $P_{24}=616\equiv 8\pmod{16}$, and $P_{25}=816\equiv 0\pmod{16}$, 
	we obtain
	\begin{align*}
		P_{28(k+1)+8} = P_{(28k+10)+26}
		&= P_{25}P_{28k+9}+P_{24}P_{28k+10}
		+P_{23}P_{28k+8}\\
		&\equiv 8P_{28k+10}+7\pmod{16}.
	\end{align*}
	Since $28k+10\equiv 3\pmod{7}$, Lemma~\ref{L2} gives 
	$\nu_2(P_{28k+10})=\nu_2(28k+14)+1=\nu_2(14(2k+1))+1=2$. Therefore, $P_{28k+10}=4q$ for 
	some odd $q$, and 
	$8P_{28k+10}=32q\equiv 0\pmod{16}$. Thus,
	\begin{equation*}
		P_{28(k+1)+8}\equiv 7\pmod{16}.
	\end{equation*}
	
	\medskip
	\noindent\textbf{Case 5.} $n\equiv 22\pmod{56}$.
	
	We will prove by induction that $P_n\equiv 31\pmod{64}$, so
	$\nu_2(P_n+1)=\nu_{2}(64t+32)=\nu_{2}(32(2t+1))=5$, where $t$ is a non-negative integer. Let $n=56k+22$ for an integer $k\ge 0$. Since $P_{22}=351\equiv 31\pmod{64}$,  the base case holds. Suppose that 
	$P_{56k+22}\equiv 31\pmod{64}$ for some fixed $k\ge 0$. Applying Lemma~\ref{L1} 
	with $r=56k+24$ and $s=54$, and using
	$P_{51}=1221537\equiv 33\pmod{64}$, $P_{52}=1618192\equiv 16\pmod{64}$, and $P_{53}=2143648\equiv 32\pmod{64}$, we obtain
	\begin{align*}
		P_{56(k+1)+22} = P_{(56k+24)+54}
		&= P_{53}P_{56k+23}+P_{52}P_{56k+24}
		+P_{51}P_{56k+22}\\
		&\equiv 32P_{56k+23}+16P_{56k+24}
		-1\pmod{64}.
	\end{align*}
Since $56k+23\equiv 2\pmod{7}$, Lemma~\ref{L2} gives that $P_{56k+23}$ is odd. Thus, $32P_{56k+23}\equiv 32\pmod{64}$. Since $56k+24\equiv 3\pmod{7}$, Lemma~\ref{L2} gives $\nu_2(P_{56k+24})=\nu_2(56k+28)+1=\nu_2(28(2k+1))+1=3$. Thus, $P_{56k+24}=8w$ for some odd $w$, and $16P_{56k+24}\equiv 0\pmod{64}$. Therefore,
	\begin{equation*}
		P_{56(k+1)+22}\equiv 31\pmod{64}.
	\end{equation*}
	
	\medskip
	\noindent\textbf{Case 6.} $n\equiv 106\pmod{112}$.
	
Let $n=112s+106$ for some integer $s\ge 0$. Then, $n+6=112(s+1)$. If $s\ge 1$, then, by the fundamental theorem of arithmetic, $s+1$ can be written uniquely as $2^{t}k$ for some integer $t\ge 0$ and some odd $k$. Thus, every $n$ in this class of residues can be uniquely expressed as $n=7\cdot 2^{t+4}k-6$ for some $t\ge 0$ and some odd $k$. First, from Lemma \ref{eq:binet_padovan} and the fact that $P_{-6}=-1$ (by Lemma \ref{eq:binet_padovan}), we have the identity
	\begin{equation}\label{eq:key-case7}
		P_n+1 = P_n - P_{-6} 
		= c_{\alpha}\alpha^{-6}(\alpha^{n+6}-1)
		+c_{\beta}\beta^{-6}(\beta^{n+6}-1)
		+c_{\gamma}\gamma^{-6}(\gamma^{n+6}-1).
	\end{equation}
	We now determine the 2-adic structure of 
	$\alpha^{n+6}-1=\alpha^{7\cdot 2^{t+4}k}-1$ 
	by iteratively squaring. For any polynomial 
	$Q(X)\in\mathbb{Z}[X]$, we compute $Q(\alpha)$ 
	by taking the remainder of the division by the 
	minimal polynomial $\phi(X)$. In this way, 
	we obtain the following chain of congruences:
	\begin{eqnarray}
	\label{eq:alpha}
		\alpha^7 
		&=& 2\alpha^2+2\alpha+1
		= 1+2(\alpha^{2}+\alpha),\nonumber\\
		\alpha^{7\cdot 2} 
		&=& 1+2^2(3\alpha^2+4\alpha+2),\\
		\alpha^{7\cdot 2^2} 
		&=& 1+2^3(77\alpha^2+102\alpha+58),\nonumber\\
		\alpha^{7\cdot 2^3} 
		&=& 1+2^4(101137\alpha^2+133978\alpha+76346),\nonumber\\
		\alpha^{7\cdot 2^4} 
		&=& 1+2^5 f(\alpha),\nonumber
	\end{eqnarray}
	where
	\begin{equation*}
		f(\alpha) := 348972965593\alpha^2
		+462290754114\alpha+263431897850.
	\end{equation*}
	Each step involves squaring the previous expression 
	and reducing it modulo $\phi(\alpha)=0$.
	
	Now, by raising $\alpha^{7\cdot 2^4}=1+2^5 f(\alpha)$ 
	to the power $2^tk$ using the binomial theorem, we obtain
	\begin{equation*}
		\alpha^{7\cdot 2^{t+4}k} 
		= \sum_{s=0}^{2^t k}\binom{2^t k}{s} 2^{5s}f(\alpha)^s.
	\end{equation*}
	For $s\ge 2$, Kummer's theorem gives 
	$\nu_2\binom{2^t}{s}=t-\nu_2(s)$, and since 
	$\nu_2(s)\le \log s/\log 2\le 5s-9$ for $s\ge 2$ 
	(since the function $h(s):=5s-9-\log s/\log 2$ vanishes 
	at $s=2$ and has a positive derivative $5-1/(s\log 2)>0$ 
	for all $s\ge 2$), 
	we obtain
	\begin{equation*}
		\nu_2\left(\binom{2^t}{s}2^{5s}\right) 
		\ge t-\nu_2(s)+5s \ge t+9\qquad {\text{\rm for}}\qquad s\ge 2.
	\end{equation*}
	Thus, 
	$$\alpha^{7\cdot 2^{t+4}}\equiv 1+2^{t+5}f(\alpha)\pmod {2^{t+9}}.
	$$
Since $k$ is odd, raising to the power $2^t k$ instead of $2^t$ yields the same congruence modulo $2^{t+9}$, since the additional factor of $k$ merely multiplies the leading term by the odd integer $k$ and leaves the higher-order terms still divisible by $2^{t+9}$. Therefore,
	\begin{equation}\label{eq:alpha-congruence}
		\alpha^{7\cdot 2^{t+4}k} 
		\equiv 1+2^{t+5}kf(\alpha)\pmod{2^{t+9}},
	\end{equation}
	and similarly for $\beta$ and $\gamma$. Substituting 
	\eqref{eq:alpha-congruence} into \eqref{eq:key-case7} 
	and using Lemma \ref{eq:binet_padovan} gives
	$$
		P_n+1 
		\equiv 2^{t+5}k\bigl(c_{\alpha}\alpha^{-6}f(\alpha)
		+c_{\beta}\beta^{-6}f(\beta)
		+c_{\gamma}\gamma^{-6}f(\gamma)\bigr)\equiv 2^{t+5}k\,M\pmod{2^{t+9}},
	$$
	where
	$$
		M := 348972965593P_{-4}
		+462290754114P_{-5}
		+263431897850P_{-6}.
	$$
	Here, we used the identity 
	$$c_{\alpha}\alpha^{-6}f(\alpha)+c_{\beta}\beta^{-6}f(\beta)
	+c_{\gamma}\gamma^{-6}f(\gamma)=348972965593P_{-4}+462290754114P_{-5}	+263431897850P_{-6},$$
 which follows from Lemma \ref{eq:binet_padovan}. Since $k$ is odd and 
 $$
 \nu_2(M)=\nu_2(198858856264)=\nu_{2}(2^3\cdot 1229\cdot 20225677)=3,
 $$ 
 we conclude that
	\begin{equation*}
		\nu_2(P_n+1) = t+5+\nu_2(k)+\nu_2(M) 
		= t+8.
	\end{equation*}
	Since $\nu_2(n+6)=\nu_2(7\cdot 2^{t+4}k)=t+4$, this gives
	\begin{equation*}
		\nu_2(P_n+1) = \nu_2(n+6)+4.
	\end{equation*}
	
	\medskip
	\noindent\textbf{An exceptional case.} $n\equiv 50\pmod{112}$.
	
	Unlike the previous cases, the residue class $n\equiv 50\pmod{112}$ does not admit 
	a uniform closed-form expression for $\nu_2(P_n+1)$. 
	Instead, we determine the maximum possible value of 
	$\nu_2(P_n+1)$ over all $n$ in this class with 
	$n<2.3\times 10^{15}$ using the following algorithm, which 
	is essentially based on Lemmas~\ref{lem:doubling} 
	and~\ref{lem:shift}.
	
	\medskip
	\noindent\textit{Algorithm.} Set $n_0=50$ and note that $\nu_2(P_{50}+1)=\nu_2(922112)=\nu_{2}(2^{9}\cdot 1801)=9$. At each step, given a current index 
	$n_r$, we compute $\nu_2(P_{n_r+7\cdot 2^{r+1}}+1)$ 
	using Lemma~\ref{lem:shift} applied with 
	$U_n=P_n$ modulo $2^{100}$, after first computing 
	$R_{7\cdot 2^k}\pmod{2^{100}}$ and 
	$R_{-7\cdot 2^k}\pmod{2^{100}}$ for 
	$k=0,1,\ldots,100$ using the doubling recurrence relations from 
	Lemma \ref{lem:doubling}. If adding $7\cdot 2^{r+1}$ 
	to the current index increases the value of $\nu_2(P_n+1)$, we set 
	$n_{r+1}=n_r+7\cdot 2^{r+1}$; otherwise we try 
	$7\cdot 2^{r+2}$, and so on. The algorithm terminates 
	when the next step would require an index greater than 
	$2.3\times 10^{15}$, which occurs approximately when
	$r=49$, since $7\cdot 2^{49}>2.3\times 10^{15}$. 
	The maximum value of $\nu_2(P_n+1)$ produced by this 
	algorithm is the constant $k_1$.
	
An implementation of this algorithm in SageMath yields 
\begin{equation*}
	k_{1}=\nu_{2}(P_{919832845308882}+1)= 56.
\end{equation*}
See Appendix \ref{app1} for more information. This completes the proof of Proposition \ref{prop:Pn}.
\end{proof}

\end{document}
